% Part: sets-functions-relations
% Chapter: sets
% Section: unions-and-intersections

\documentclass[../../../include/open-logic-section]{subfiles}

\begin{document}

\olfileid{sfr}{set}{uni}
\olsection{সংযোগ ও ছেদ}

\begin{explain}
\olref[sfr][set][bas]{sec}-এ ধর্মনির্দেশের সাহায্যে সেটের সংজ্ঞা,
অর্থাৎ $\Setabs{x}{\phi(x)}$ আকারের সংজ্ঞা প্রবর্তন করেছি। এখানে
কোনও ধর্ম~$\phi$ ব্যবহার করা হয়; এই ধর্মে আগে সংজ্ঞায়িত সেটের উল্লেখ
থাকতে পারে। যেমন, $A$ ও~$B$ সেট হলে,
$\Setabs{x}{x \in A \lor x \in B}$ সেটটি সেই সমস্ত বস্তু নিয়ে গঠিত,
যেগুলি $A$ অথবা~$B$-এর !!{element}s। অর্থাৎ, এই সেটে $A$ ও~$B$-এর
!!{element}s একত্র করা হয়। \olref{fig:union}-এর মতো করে এটি
দৃশ্যায়িত করা যায়: চিহ্নিত অঞ্চলটি $A$ ও~$B$, দুই সেটের
!!{element}s-কে একত্রে নির্দেশ করে।

\begin{figure}
  \olasset{assets/diagrams/union.tikz}
  \caption{দুই সেটের সংযোগ $A \cup B$ হল $A$ ও~$B$-এর
    !!{element}s একত্র করে গঠিত সেট।}
  \ollabel{fig:union}
\end{figure}

সেটগুলিকে একত্র করার এই ক্রিয়াটি খুবই উপযোগী ও প্রচলিত।
তাই এর একটি নির্দিষ্ট নাম ও প্রতীক দেব।
\end{explain}

\begin{defn}[সংযোগ]
দুটি সেট $A$ ও $B$-এর \emph{সংযোগ}, যাকে $A \cup B$ লেখা হয়,
হল সেই সমস্ত বস্তুর সেট, যেগুলি $A$, $B$, অথবা উভয়ের !!{element}s।
\[
A \cup B = \Setabs{x}{x \in A \lor x \in B}
\]
\end{defn}

\begin{ex}
!!{element}s কত বার আছে, তা বিবেচ্য নয়। তাই দুটি সেটে সাধারণ
!!a{element} থাকলে, তাদের সংযোগে সেই !!{element} এক বারই থাকে।
যেমন, $\{ a, b, c\} \cup \{ a, 0, 1\} = \{a, b, c, 0, 1\}$।

কোনও সেট এবং তার একটি উপসেটের সংযোগ হল বড় সেটটিই:
$\{a, b, c \} \cup \{a \} = \{a, b, c\}$।

কোনও সেটের সঙ্গে শূন্য সেটের সংযোগ সেই সেটেরই সমান:
$\{a, b, c \} \cup \emptyset = \{a, b, c \}$।
\end{ex}

\begin{prob}
প্রমাণ করো যে, $A \subseteq B$ হলে $A \cup B = B$।
\end{prob}

\begin{explain}
সংযোগের একটি ``দ্বৈত'' ক্রিয়াও বিবেচনা করতে পারি। এতে সেই সমস্ত
!!{element}sের সেট গঠিত হয়, যেগুলি $A$-এর !!{element}s এবং
$B$-এরও !!{element}s। এই ক্রিয়ার নাম \emph{ছেদ};
\olref{fig:intersection}-এর মতো করে এটি দেখানো যায়।
\begin{figure}
  \olasset{assets/diagrams/intersection.tikz}
  \caption{দুই সেটের ছেদ $A \cap B$ হল তাদের সাধারণ
    !!{element}sের সেট।}
  \ollabel{fig:intersection}
\end{figure}
\end{explain}

\begin{defn}[ছেদ]
দুটি সেট $A$ ও $B$-এর \emph{ছেদ}, যাকে $A \cap B$ লেখা হয়,
হল সেই সমস্ত বস্তুর সেট, যেগুলি $A$ ও~$B$, উভয়ের !!{element}s।
\[
A \cap B = \Setabs{x}{x \in A \land x \in B}
\]
দুটি সেটের ছেদ শূন্য হলে তাদের \emph{বিচ্ছিন্ন} বলে।
এর অর্থ, তাদের কোনও সাধারণ !!{element}s নেই।
\end{defn}

\begin{ex}
দুটি সেটে কোনও সাধারণ !!{element}s না থাকলে তাদের ছেদ শূন্য:
$\{ a, b, c\} \cap \{ 0, 1\} = \emptyset$।

দুটি সেটে সাধারণ !!{element}s থাকলে, সেই সব উপাদানের সেটই তাদের ছেদ:
$\{a, b, c \} \cap \{a, b, d \} = \{a, b\}$।

কোনও সেট এবং তার একটি উপসেটের ছেদ হল ছোট সেটটিই:
$\{a, b, c\} \cap \{a, b\} = \{a, b\}$।

যে-কোনও সেটের সঙ্গে শূন্য সেটের ছেদ শূন্য:
$\{a, b, c \} \cap \emptyset = \emptyset$।
\end{ex}

\begin{prob}
কঠোরভাবে প্রমাণ করো যে, $A \subseteq B$ হলে $A \cap B = A$।
\end{prob}

\begin{explain}
দুটির বেশিসংখ্যক সেটেরও সংযোগ বা ছেদ গঠন করা যায়। সাধারণভাবে
তা করার একটি সুন্দর উপায় হল: যে সেটগুলির সংযোগ বা ছেদ গঠন করতে
চাই, সেগুলিকে একটি সেটে একত্র করি। তখন মূল সেটগুলির সংযোগকে
সংজ্ঞায়িত করতে পারি সেই সমস্ত বস্তুর সেট হিসেবে, যেগুলি একত্র করা
সেটটির অন্তত একটি !!{element}ের সদস্য। আর ছেদ হবে সেই সমস্ত
বস্তুর সেট, যেগুলি ওই সেটটির প্রতিটি !!{element}ের সদস্য।
\end{explain}

\begin{defn}
$A$ সেটের একটি সেট হলে, $\bigcup A$ হল $A$-এর
!!{element}sের !!{element}sের সেট:
\begin{align*}
\bigcup A & = \Setabs{x}{x \text{ যে সেটের সদস্য, সেটি } A\text{-এর !!a{element}}},
\text{ অর্থাৎ,}\\
& = \Setabs{x}{\text{এমন একটি } B \in A
  \text{ আছে যাতে } x \in B}
\end{align*}
\end{defn}

\begin{defn}
$A$ সেটের একটি সেট হলে, $\bigcap A$ হল সেই সমস্ত বস্তুর সেট,
যেগুলি $A$-এর সমস্ত উপাদানের সাধারণ সদস্য:
\begin{align*}
\bigcap A & = \Setabs{x}{x \text{ সদস্য হয় } A\text{-এর প্রতিটি !!{element}ের}},
\text{ অর্থাৎ,}\\
 & = \Setabs{x}{\text{সমস্ত } B \in A\text{-এর জন্য}, x \in B}
\end{align*}
\end{defn}

\begin{ex}
ধরি, $A = \{ \{ a, b \}, \{ a, d, e \}, \{ a, d \} \}$।
তাহলে $\bigcup A = \{ a, b, d, e \}$ এবং $\bigcap A = \{ a \}$।
\end{ex}
\begin{prob}
	দেখাও যে, $A$ সেট হলে এবং $A \in B$ হলে $A \subseteq \bigcup B$।
\end{prob}

সেটের অনুক্রম $A_1$, $A_2$, \dots-এর ক্ষেত্রেও একই কাজ করতে পারি:
\begin{align*}
\bigcup_i A_i & = \Setabs{x}{x \text{ সদস্য হয় কোনও একটি } A_i\text{-এর}}\\
\bigcap_i A_i & = \Setabs{x}{x \text{ সদস্য হয় প্রতিটি } A_i\text{-এর}}.
\end{align*}

সেটগুলির একটি \emph{সূচক} থাকলে, অর্থাৎ এমন কোনও সেট $I$ থাকলে,
যার প্রতিটি $i \in I$-এর জন্য আমরা $A_i$ বিবেচনা করছি,
নিচের সংক্ষিপ্ত প্রতীকগুলিও ব্যবহার করতে পারি:
\begin{align*}
	\bigcup_{i \in I} A_i & = \bigcup \Setabs{A_i }{i \in I}\\
	\bigcap_{i \in I} A_i & = \bigcap\Setabs{A_i}{i \in I}
\end{align*}

শেষে, $A$-এর সেই সমস্ত !!{element}sের সেটও বিবেচনা করতে পারি,
যেগুলি $B$-এ নেই। \olref{difference}-এর মতো করে এটি দেখানো যায়।

\begin{figure}
  \olasset{assets/diagrams/difference.tikz}
  \caption{দুই সেটের অন্তর $A \setminus B$ হল $A$-এর সেই
    !!{element}sের সেট, যেগুলি $B$-এরও !!{element}s নয়।}
  \ollabel{difference}
\end{figure}

\begin{defn}[অন্তর]
\emph{সেটের অন্তর}~$A \setminus B$ হল $A$-এর সেই সমস্ত
!!{element}sের সেট, যেগুলি $B$-এরও !!{element}s নয়। অর্থাৎ,
\[
A\setminus B = \Setabs{x}{x\in A \text{ এবং } x \notin B}.
\]
\end{defn}

\begin{prob}
	প্রমাণ করো যে, $A \subsetneq B$ হলে $B \setminus A \neq \emptyset$।
\end{prob}

\end{document}
